\title{LongRoPE: Extending LLM Context Window Beyond 2 Million Tokens}

\begin{abstract}
Expanding the context window in large language models (LLMs) is highly sought after. Nonetheless, present-day context window extensions typically do not exceed approximately 128k tokens, owing to the substantial expense of fine-tuning, the limited availability of lengthy textual data, and the problematic values generated by previously unseen token positions.

In this paper, we present LongRoPE, a novel approach that expands the context window of pre-trained large language models (LLMs) to a remarkable \textbf{2048k} tokens, requiring as few as 1,000 fine-tuning steps at training lengths up to 256k tokens, all while preserving the original performance on short context lengths. This advancement is accomplished through three main contributions: (i) we pinpoint and leverage two distinct forms of non-uniformities within positional interpolation, employing an efficient search method to optimize initialization for fine-tuning, thereby enabling an 8$\times$ increase in the context window without additional fine-tuning; (ii) we propose a stepwise extension approach, where we initially fine-tune the LLM at 256k context length and subsequently apply a secondary positional interpolation to the extended model, ultimately achieving a 2048k context window; (iii) we refine LongRoPE specifically at the 8k context window to restore and sustain the performance for shorter inputs. Comprehensive evaluations using LLaMA2 and Mistral on a variety of benchmark tasks verify the effectiveness of our framework. The LongRoPE-augmented models retain the standard architecture with only slight adjustments to the positional embedding component, allowing compatibility with the majority of prior optimizations. Code will be released at \texttt{https://github.com/microsoft/LongRoPE}

\section{Introduction}

Large Language Models (LLMs) have achieved significant advancements across a variety of applications~\cite{gpt4,llama2}, yet they remain constrained by relatively short context window sizes—for instance, LLaMA2 is limited to handling 4096 tokens~\cite{llama2}. When input extends beyond this predefined window, model effectiveness deteriorates as it encounters positional contexts outside its training distribution. Such limitations are particularly problematic in key applications, including in-context learning with extensive exemplars~\cite{Huang2023FewerIM} and the deployment of LLM-powered agents~\cite{park2023generative,madaan2023selfrefine}.

Recent studies have demonstrated that it is possible to enlarge the context window of a pre-trained LLM to approximately 128k tokens by fine-tuning on documents of increased length~\cite{longlora,pi,yarn,zhang2024soaring,liu2023scaling}. Nonetheless, scaling the context window beyond this remains challenging due to three principal factors. Firstly, introducing new, untrained positional indices yields numerous anomalous values, resulting in significant out-of-distribution complications and rendering the fine-tuning process hard to stabilize~\cite{pi}. This issue is exacerbated when context extension is performed from 4k to $>$1000k, as this process introduces over 90\% unseen positions. Secondly, effective fine-tuning typically depends on access to training samples whose lengths match those of the target context. However, datasets containing sufficiently long sequences—particularly those over 1000k tokens—are scarce. In addition, training on such ultra-long sequences incurs extremely high computational costs, necessitating vast training durations and considerable GPU capacity. Thirdly, as the context window is drastically increased, the attention mechanism must distribute focus across a substantially larger set of token positions. This dilution weakens model performance on tasks involving short contexts, as attention becomes overly diffuse~\cite{pi}.

A commonly adopted strategy to address the first issue involves interpolating RoPE positional embeddings~\cite{rope,pi}, which rescales novel positional indices so they fall within the pretrained domain, as illustrated in Fig.~\ref{fig:overview}. In particular, Position Interpolation (PI)~\cite{pi} linearly adjusts RoPE’s rotary angles by the extension proportion. In contrast, NTK~\cite{ntk,dynamicntk} introduces non-uniform interpolation and extrapolation across the various RoPE dimensions. YaRN~\cite{yarn} partitions RoPE dimensions into three groups according to frequency, subsequently assigning each group to extrapolation, NTK, or linear interpolation as appropriate. Despite these innovations, positional embedding within the Transformer model demonstrates a \emph{complex, non-uniform information entropy}. Current methodologies fail to fully exploit these intricate entropy characteristics, resulting in information degradation and thereby restricting the feasible context window length.

Section~\ref{sec:analysis} empirically demonstrates two main insights: \textbf{(1)} Successful positional interpolation requires addressing two sources of non-uniformity—specifically, the variability across RoPE dimensions and token positions. Lesser interpolation is advantageous for lower RoPE dimensions and tokens appearing earlier in the sequence, yet the best choices vary depending on the desired extended context length. 
\textbf{(2)} Integrating these non-uniformities into the interpolation process preserves more of the original RoPE’s information, particularly in crucial dimensions and for important token positions. As a result, the adverse effects of interpolation are reduced, leading to improved initial parameters for subsequent fine-tuning. Additionally, this approach supports up to an 8$\times$ increase in length even without additional fine-tuning.

Inspired by these results, we present \textbf{LongRoPE}, a robust technique designed to stretch the LLM context window to lengths exceeding 2 \emph{million} tokens. The development of LongRoPE rests on three fundamental advancements. Firstly, {LongRoPE} leverages multidimensional non-uniformities inherent in positional interpolation to their fullest. It discovers optimal rescale factors for the rotational angles in RoPE, tailoring them for each individual RoPE dimension according to the tokens’ positions. Given that the process of identifying these rescale factors becomes exponentially complex as the desired extension increases, {LongRoPE} incorporates an evolutionary search strategy enhanced by two specific optimization methods to improve the efficiency of the search. An illustration of the resulting rescaled RoPE can be seen in Fig.~\ref{fig:overview}.

Subsequently, {LongRoPE} utilizes a resourceful, gradual extension approach to reach a 2048k context window, circumventing the necessity for explicit fine-tuning on ultra-long documents, which are both rare and not readily accessible. The process initiates with identifying a 256k token length compatible with the pre-trained LLM, followed by fine-tuning the model at this sequence length. With our non-uniform positional interpolation enabling an extension by a factor of 8$\times$ even when skipping further fine-tuning, a second optimization is performed to determine suitable new RoPE rescale factors on the already fine-tuned, length-extended model. This procedure successfully enables a 2048k context window for both LLaMA2 and Mistral~\cite{mistral}.

Lastly, to prevent reduced performance on the original, shorter context windows, {LongRoPE} maintains tuning of the RoPE rescale parameters in the expanded LLM. In parallel with the upscaling from 256k to 2048k, we utilize our search algorithm to downscale to 4k and 8k context windows on the 256k fine-tuned LLM, aiming to minimize positional interpolation. For inference, when the input sequence is shorter than 8k, RoPE is modified with the rescale factors identified by our search process.

A comprehensive set of experiments involving multiple LLMs and a range of long-context benchmarks confirms the efficacy of our approach. Our findings indicate that LongRoPE consistently preserves low perplexity across evaluation lengths spanning from 4k up to 2048k, attains greater than 90\% passkey retrieval accuracy, and achieves accuracy levels on par with existing models when assessed on standard benchmarks with a 4096 token context window. The LongRoPE technique is compatible with all LLMs utilizing RoPE embeddings. We intend to make both our code and the LongRoPE-2048k models publicly available.

\section{Non-uniformity in Positional Interpolation}
\label{sec:analysis}

\subsection{Preliminary}

Transformer models require explicit positional information, often in the form of position embedding, to represent the order of input tokens. Our work focuses on the RoPE~\cite{rope} position embedding, which is widely used in recent  LLMs. For a token at position index $n$, its corresponding RoPE encoding can be simplified as follows:

\begin{equation}
		\fontsize{7.8}{7.8} \selectfont
	\label{eq:rope}
	\left[ cos(n\theta_0), sin(n\theta_0), cos(n\theta_1), \cdot\cdot\cdot, cos(n\theta_{d/2-1}), sin(n\theta_{d/2-1}) \right]   
\end{equation}

Here, $d$ denotes the size of the embedding, $n\theta_i$ indicates the rotational angle associated with the token located at position $n$, and $\theta_i=\theta^{-2i/d}$ encodes the frequencies used for rotation. In the RoPE mechanism, the typical base value assigned to $\theta$ is 10000.

\textbf{\emph{Context window extension ratio $s$} and positional interpolation}. The ratio $s$ is specified as the quotient of the new, extended context length $L'$ to the original context length $L$, i.e., $s=\frac{L'}{L}$.

To extend the context window from $L$ to $L'$, current positional interpolation methods suggest downscaling rotation frequencies $\theta_i$ based on extension ratio $s$. Let $\beta=\theta^{2/d}$, and $\lambda$ denote the actual rescale factor related to $s$, we   unify these positional interpolation methods as follows:

\begin{equation}	
	\fontsize{7.5}{7.5} \selectfont
	\label{eq:unified_rope}
	\begin{aligned}
		\left[cos\left(\frac{n}{\lambda(\beta)^0}\right), sin \left(\frac{n}{\lambda(\beta)^0}\right), cos\left(\frac{n}{\lambda(\beta)^1}\right), \cdot\cdot\cdot,  sin \left(\frac{n}{\lambda(\beta)^{d/2-1}}\right) \right]   
	\end{aligned}
\end{equation}

\textbf{Linear positional interpolation (PI)}. PI~\cite{pi} introduces a technique where position indices are interpolated linearly, confined to the range of sequence lengths encountered during pre-training. When scaling to a sequence length determined by an extension ratio $s$, all position encoding rotation angles are uniformly contracted by a factor of $\lambda=s$ across every RoPE dimension. This uniform compression causes position encodings to become densely clustered, making it challenging for the model to effectively separate tokens that are close to each other in the sequence. Consequently, PI typically shows suboptimal performance when applied to scenarios involving large extension ratios.

\textbf{NTK-based interpolation and extrapolation}. The works of~\cite{ntk,dynamicntk} approach RoPE through the lens of information encoding and utilize insights from Neural Tangent Kernel (NTK) theory~\cite{jacot2018neural,tancik2020fourier}. To address the problem of positional congestion in PI, their method disperses the effects of interpolation across the various RoPE dimensions. This is achieved by applying a smaller scaling factor to the lower (high-frequency) dimensions and a larger scaling to the higher (low-frequency) ones, thereby enabling both positional interpolation and extrapolation, formalized as $\lambda=s^{i}$. The refined dynamic NTK approach~\cite{dynamicntk} further customizes the extension ratio at each positional index depending on the sequence's current length. In contrast to PI—which requires additional fine-tuning—NTK-based solutions allow for context window extension without the need for fine-tuning, though the maximal achievable extension ratio is typically limited to 4$\times$.

\textbf{YaRN}~\cite{yarn} presents a notable advance in enhancing positional interpolation. The method segments the RoPE dimensions into three distinct categories according to their frequencies, assigning a tailored interpolation technique to each group. Extrapolation ($\lambda$=1) is applied to high-frequency dimensions, while linear interpolation (PI) is reserved for those with low frequency, and intermediate frequencies are handled by the NTK approach. The principal contribution of YaRN is this frequency-based segmentation of RoPE dimensions. However, the determination of these groupings currently relies on manual, empirically-driven approaches, potentially leading to sub-optimal results when applied to newly developed LLM architectures.

\subsection{Study on Non-uniform Positional Interpolation}

Drawing motivation from NTK and YaRN, we observe that their improvements stem largely from the incorporation of non-linear elements, especially through treating various RoPE dimensions according to their frequency characteristics to enable targeted interpolation and extrapolation. Nonetheless, prevailing approaches to non-linearity are predominantly dependent on heuristics crafted by humans. Consequently, this prompts two central inquiries: (1) Does the present method for positional interpolation represent the best possible approach? (2) Are there additional forms of non-linearity that remain to be discovered?

To investigate these issues, we employ evolution search (refer to Sec.~\ref{sec:method}) to identify improved non-uniform positional interpolations tailored for LLaMA2-7B. Perplexity serves as the search criterion, evaluated over 5 randomly selected instances from the PG19~\cite{pg19} validation dataset. Our empirical evaluation leads to several notable observations.

\textbf{Finding 1}: \textit{RoPE dimensions exhibit substantial non-uniformities, which are not effectively handled by current positional interpolation methods.}

We perform a search for the optimal $\lambda$ corresponding to each RoPE dimension as specified in Eq.~\ref{eq:unified_rope}. Table~\ref{tbl:exp1} presents a comparison of the perplexity scores for LLaMA2-7B, across various approaches, evaluated on the PG19 and Proof-pile~\cite{proof-pile} test datasets, all conducted without additional fine-tuning. The results from our search yield notable improvements, indicating that both the conventional linear (PI) and the more recent non-uniform (Dynamic-NTK and YaRN) interpolation methods are not optimal. Interestingly, on the PG19 benchmark, YaRN performs worse than both PI and NTK, primarily because, in the absence of further fine-tuning, it fails to span the intended context window size. Specifically, YaRN exhibits a sharp rise in perplexity beyond 7k tokens within an 8k context window.

In our exploration, we find that the rescaled factors $\lambda$ in Eq.~\ref{eq:unified_rope} exhibit non-uniformity, diverging from the constant scaling factor $s$ used in PI, NTK's formula, and YaRN's group-wise approach. Utilizing these variable scaling factors leads to notable enhancements in the language modeling capabilities (as measured by perplexity) of LLaMA2 for extended context lengths of 8k and 16k, all achieved without additional fine-tuning. This improvement can be attributed to the newly designed positional embedding's ability to more faithfully maintain the structure of the original RoPE—particularly in the principal dimensions—thereby making it easier for the LLM to differentiate between tokens in nearby positions.

\textbf{Finding 2}: \textit{RoPE for the initial tokens in the input sequence should be extrapolated with less interpolation.} 

We propose that for the first $\hat{n}$ tokens in the input, the application of RoPE-based interpolation should be minimized. The rationale behind this stems from their disproportionately high attention scores, underscoring their importance within attention mechanisms, a trend also identified by Streaming LLM~\cite{streamingllm} and LM-Infinite~\cite{han2023lminfinite}. To assess this, we experimented with expanding the context window to lengths of 8k and 16k using PI and NTK approaches, systematically varying the count of non-interpolated initial tokens $\hat{n}$ across values (0, 2, ..., 256). Setting $\hat{n}$ to 0 restores the standard PI and NTK procedures. As shown in Table~\ref{tbl:tokenposition}, two notable findings emerge: \textbf{(1)} exempting the initial tokens from position interpolation leads to notable improvements in model performance; \textbf{(2)} the best choice of $\hat{n}$ is contingent upon the desired extent of the context window expansion.

\textbf{Finding 3}: \textit{Non-uniform positional interpolation effectively extends LLM context window in both fine-tuning and non-fine-tuning settings.} 

Although our findings demonstrate that leveraging non-uniform position interpolation markedly enhances extension capability at 8k and 16k context lengths without additional fine-tuning, scaling to even longer contexts necessitates further fine-tuning. To this end, we adapt LLaMA2-7B using our optimized RoPE approach, extending the context window to 64k tokens (refer to the Appendix for detailed configuration). As indicated in Table~\ref{tbl:finetune}, this approach yields considerably stronger results compared to both PI and YaRN, regardless of whether fine-tuning has been applied to LLaMA2-7B. This improvement stems from the superior handling of non-uniform positional interpolation, which mitigates information degradation and establishes more effective fine-tuning initial conditions.

\textbf{Summary}. Our research identifies two sources of non-uniformity—across both RoPE dimensions and token indices. Effectively incorporating these aspects into positional interpolation methods substantially boosts the performance of large language model context extensions.

\section{LongRoPE}

\label{sec:method}
Building on these insights, we propose LongRoPE, an approach that initially incorporates an effective search algorithm to leverage both forms of non-uniformity. This method is subsequently employed to expand the LLM context window to surpass 2 million tokens.

\subsection{Problem Formulation}

These two forms of non-uniformity contribute to an expansive solution space and complicate the optimization process. To manage this, we recast the multidimensional non-uniform position interpolation optimization issue as a search-based problem.

For a LLM targeting a  context window size of $L'$ and lengthy input documents $\mathbf{X}$, where each $\mathbf{x}\in\mathbf{X}$ surpasses $L'$ in token length, we denote the original rotary angle of the $i^{th}$ dimension in RoPE embedding at token position $n$ as  $\frac{n}{\beta^i}$. The optimization problem is then formulated as follows:

\begin{equation}
\begin{gathered}
		\label{eq:problem}

\underset {\mathbf{x}\in\mathbf{X}; \, |\mathbf{x}|\ge L'}{ \operatorname {arg\,min} } \,  \mathcal{L}\, \left( \text{LLM} (\text{RoPE}, \mathbf{X})\right),	\text{where \;} 
\\
\scalebox{0.93}{
$\underset {\begin{subarray}{l} i=0,\cdot\cdot,\frac{d}{2}-1;\\ n\in[ 0,|\mathbf{x}|);\end{subarray}}{\text{RoPE}_(n)}= {\left[\cdot\cdot,\cos\left(\mathbb{I}(\hat{\lambda}_{i}, \hat{n})\frac{n}{\beta^i}\right), \sin\left(\mathbb{I}(\hat{\lambda}_{i}, \hat{n})\frac{n}{\beta^i}\right),\cdot\cdot\right] }$}\\
	\text{where \;} \mathbb{I}(\hat{\lambda}_{i},\hat{n})=\begin{cases}
	1&\mbox{\;  $n<\hat{n}$}\\
{\frac{1}{\lambda}_{i}}&\mbox{\; $n\geq\hat{n}$}
\end{cases}
	\end{gathered}
\end{equation}

In this approach, we define a set of scaling factors, $\mathbb{I}(\hat{\lambda}_{i},\hat{n})$, to address two types of non-uniformities. Here, $\hat{\lambda_i}$ characterizes the degree of non-uniformity across RoPE dimensions, while $\hat{n}$ specifies the variation in token positions. The scaling factor $\mathbb{I}(\hat{\lambda}_{i},\hat{n})$ is employed to adjust the rotation angle associated with the $i^{th}$ RoPE dimension, with $\hat\lambda_{i}$ serving as the specific scaling coefficient and $\hat{n}$ functioning as a threshold on token position. For the first $\hat{n}-1$ token indices, the original rotary position encoding angle $\frac{n}{\beta^i}$ remains unchanged, and $\hat\lambda_{i}$ is not applied. Conversely, once the token position $n$ satisfies $n \ge \hat{n}$, the rescaling factor becomes active.

Assuming a desired context window length of $L'$, our aim is to determine the best rescale factors ($\mathbb{I}(\hat{\lambda}_{0},\hat{n})$, $\mathbb{I}(\hat{\lambda}_{1},\hat{n})$, ... $\mathbb{I}(\hat{\lambda}_{i},\hat{n})$, ...) for each of the RoPE dimensions ranging from the $1^{st}$ to the $d^{th}$. Through the application of these optimized rescale values, the target $\text{LLM}$ with its modified $\text{RoPE}$ is expected to attain the lowest possible next token prediction loss, $\mathcal{L}$ (specifically, perplexity) on inputs $\mathbf{X}$ where the number of tokens equals $L'$.

\subsection{Searching the Non-uniform Position Interpolation}

To address the challenge outlined in Eq.~\ref{eq:problem}, we propose an efficient and straightforward technique that identifies the best RoPE rescaling parameters, thereby maximizing the utility of multidimensional irregularities inherent in position embedding.

\textbf{Search space}. We construct an extensive search space that incorporates both forms of non-uniformity. The structure of this search space is summarized in Table~\ref{tbl:searchspace}. Rather than assigning a uniform rescale factor across all dimensions, our approach enables the selection of unique rescale factors for each dimension in the RoPE embedding. For the sake of simplifying the search space, we opt to search over $\lambda_i$ and $\hat{n}$ directly, instead of searching the set $\mathbb{I}(\hat{\lambda}_{i},\hat{n})$, where we have $\hat{\lambda}_{i} = 1/\lambda_i$. According to Table~\ref{tbl:searchspace}, the permissible range for each $\lambda_i$ starts at 1.0 (corresponding to straightforward extrapolation) and extends up to $s\times1.25$ (indicating broader interpolation than that used in PI), with increments of 0.01. Here, $s$ denotes the ratio by which the target context window is extended.

$\hat{n}$ determines how many of the initial token positions preserve the original RoPE embedding, foregoing position interpolation. In practice, we select $\hat{n}$ from the set \{0, 1, 2, 4, 8, 12, 16, 20, 24, 28, 32, 64, 128, 256\} during search. If $\hat{n}=0$, then every token position applies the optimized rescale factors.

\textbf{Evolution-based search}. The search space presented in Table~\ref{tbl:searchspace} encompasses a vast array of potential positional interpolation configurations, making effective navigation of this space highly complex. As an illustration, selecting a $s=4\times$ extension produces $400^{128/2}\times14$=4$\times10^{167}$ possible options. Notably, as the extension ratio increases, the search space grows exponentially. To manage this complexity, we adopt evolution search~\cite{guo2020single} and implement two optimization strategies that substantially improve search efficiency. The overall procedure is outlined in Algorithm~\ref{alg:search}.

\textit{Optimized initial population generation}. Rather than selecting all $P$ rescale factors at random to form the initial population, we explicitly include the RoPE rescale factors associated with PI, NTK, and YaRN as part of the initial individuals. The remaining $P$-3 members of the population are derived by applying random mutations to these three rescale factors, each mutation occurring with probability $p$.

\textit{Monotonically non-decreasing constraint}. Once the initial population is produced, we evaluate LLM perplexity for every candidate. This involves modifying the target LLM with the relevant RoPE rescale factors and assessing the perplexity on input $\mathbf{X}$. The best-performing k individuals are then selected as parents for subsequent generations. Nevertheless, due to the expansive search space, straightforward mutation and crossover operations may frequently produce suboptimal solutions, resulting in excessive perplexity evaluations. This inefficiency becomes especially pronounced for large $L'$, since each perplexity assessment requires substantial inference time.

To this end, we enforce a monotonicity condition such that the rescaled RoPE factors adhere to $\lambda_i \le \lambda_{i+1}$. Only those RoPE parameterizations meeting this non-decreasing requirement are utilized in LLM perplexity assessments, which substantially curtails the scope of the parameter search. More concretely, we stipulate that $\lambda_i$ must grow monotonically with respect to the RoPE dimension index ($i=0, ..., 63$). The rationale for imposing this monotonicity along the dimensions draws on NTK theory~\cite{jacot2018neural,tancik2020fourier,ntk}, which posits that the lower-index (higher-frequency) dimensions necessitate less interpolation—meaning smaller $\lambda_i$ values—whereas higher-index (lower-frequency) dimensions can accommodate larger $\lambda_i$ values, corresponding to more interpolation.

\begin{algorithm}[t]
	\small
	\caption{The search algorithm}
	\textbf{Input:} target LLM, input samples $\mathbf{X}$,   population size $P$, mutation size $N_1$, crossover size $N_2$, max iterations $\mathcal{T}$, mutate probability $p$\\

	\begin{algorithmic}[1]
		\label{alg:search}
		\STATE Top-k$\gets\phi$;
		\STATE $\text{P}_0\gets$\textit{Initialize\_population\_with\_optimization}($P$, $\mathbf{X}$, $p$);
		\FOR{$i$ = $1$ \TO $\mathcal{T}$}
		\STATE \textit{Compute\_perplexity}(LLM, $\text{P}_{i-1}$, $\mathbf{X}$);
		\STATE Top-k $\gets$ \textit{Update\_Topk}(Top-k, $\text{P}_{i-1}$);
		\STATE $\text{P}_{mutation} \gets$ \textit{Mutation\_with\_mono\_constraint}(Top-k, $N_1$, $p$);
		\STATE $\text{P}_{crossover} \gets$ \textit{Crossover\_with\_mono\_constraint}(Top-k, $N_2$);
		\STATE $\text{P}_i \gets \text{P}_{mutation} \cup \text{P}_{crossover} \cup \text{Top-k}$;
		\ENDFOR
		\STATE Output the candidate with the minimal perplexity from Top-k;
	\end{algorithmic}
\end{algorithm}

\textbf{8$\times$ extension without fine-tuning}. By utilizing evolutionary search, our approach successfully determines non-uniform RoPE rescale factors that maintain critical positions and dimensions, thereby reducing the loss of information caused by interpolation. As illustrated in Fig.\ref{fig:non-ft}, this technique allows LLaMA2's context window to be expanded from 4k to 32k without the need for fine-tuning. In comparison, other strategies like PI, as well as non-uniform NTK and YaRN, result in significant increases in perplexity beyond a 2$\times$ extension.

\subsection{Extending LLM Context Window to 2048K}
\label{sec:extendto2048k}

\textbf{Progressive extension to 2048k}. In this section, we present our approach for scaling the context window of pre-trained LLMs beyond the standard 4k limit, reaching up to 2048k. Our non-uniform positional interpolation technique allows for an 8$\times$ expansion without any need for fine-tuning. However, to achieve significantly larger extensions, such as 512$\times$, fine-tuning becomes unavoidable. A conventional strategy involves determining optimal RoPE rescale factors corresponding to the desired 2048k length, followed by model fine-tuning. Nevertheless, this process demands extremely high computational resources, making it challenging in practice. In addition, our experience indicates that effectively fine-tuning LLMs under such extreme extension ratios is quite difficult (refer to Appendix for further discussion).

Fortunately, LongRoPE demonstrates efficacy with both pre-trained models and those that have undergone fine-tuning for extended sequence lengths. To this end, we propose a streamlined, incremental approach that reaches the desired 2048k target using only 1k fine-tuning iterations and a maximum training length of 256k.

$\Diamond$ \textit{Scaling pre-trained LLMs to 256k via LongRoPE search}. Using LLaMA2 as the baseline, we experiment with target context window lengths of 128k and 256k. At these levels, the models are extended by factors of 32$\times$ and 64$\times$, correspondingly.

$\Diamond$ \textit{Fine-tuning to 256k}. In the subsequent phase, we adapt the pre-trained LLM to support a context window of 256k. To accomplish this, we initiate fine-tuning of LLaMA2 for 400 steps employing RoPE rescaled factors set for 128k. After this stage, we switch the RoPE rescaled factors to correspond with 256k on the obtained checkpoint, and then carry out a further 600 fine-tuning steps. This staged approach offers greater efficiency compared to directly fine-tuning on the 256k context size from the start.

$\Diamond$ \textit{Applying LongRoPE search to increase context window to 2048k on fine-tuned LLM}. In the final step, we conduct an additional search process on the LLM previously fine-tuned at a 256k context length. This approach allows the model’s context window to be extended substantially to 2048k, all without the need for any more fine-tuning. Thus, the overall extension achieves a $512\times$ increase.

\textbf{Restoring performance in shorter context windows}. Expanding the context window to a substantial size of 2048k results in reduced performance within the initial context range. This phenomenon is well-documented in positional interpolation~\cite{pi}, which compresses position embeddings for the original context window into a significantly smaller subset of the embedding space, thereby degrading language model efficacy. When employing a 512$\times$ expansion factor, the representation of positions in the standard 4k window is especially densely packed.

To address this, an additional evolution search is conducted on the extended LLM, focusing on optimizing the RoPE rescale factors specifically for shorter context lengths such as 4k and 8k. The upper bound for the searched $\lambda$ is lowered because less positional interpolation is necessary at these reduced lengths. When running inference, the LLM automatically modifies the relevant RoPE rescale factors accordingly.

\section{LongRoPE}

Driven by these results, we propose LongRoPE. This approach initially implements an effective search algorithm designed to harness the two identified non-uniformities. Utilizing this mechanism, we are able to extend the context window of LLMs to exceed 2 million tokens.

\subsection{Problem Formulation}

These two types of non-uniformity significantly expand the solution space, complicating the optimization process. To manage this, we recast the multidimensional non-uniform position interpolation optimization into a search problem.

For a LLM targeting a  context window size of $L'$ and lengthy input documents $\mathbf{X}$, where each $\mathbf{x}\in\mathbf{X}$ surpasses $L'$ in token length, we denote the original rotary angle of the $i^{th}$ dimension in RoPE embedding at token position $n$ as  $\frac{n}{\beta^i}$. The optimization problem is then formulated as follows:

\begin{equation}
\begin{gathered}
		\label{eq:problem}

\underset {\mathbf{x}\in\mathbf{X}; \, |\mathbf{x}|\ge L'}{ \operatorname {arg\,min} } \,  \mathcal{L}\, \left( \text{LLM} (\text{RoPE}, \mathbf{X})\right),\quad \text{with}
\\
\scalebox{0.93}{
$\underset {\begin{subarray}{l} i=0,\ldots,\frac{d}{2}-1;\\ n\in[ 0,|\mathbf{x}|);\end{subarray}}{\text{RoPE}_(n)}= {\left[\ldots, \cos\left(\mathbb{I}(\hat{\lambda}_{i}, \hat{n})\frac{n}{\beta^i}\right), \sin\left(\mathbb{I}(\hat{\lambda}_{i}, \hat{n})\frac{n}{\beta^i}\right),\ldots\right] }$}\\
	\text{where } \mathbb{I}(\hat{\lambda}_{i},\hat{n})= \begin{cases}
	1&\mbox{ if $n<\hat{n}$}\\
	{\frac{1}{\lambda}_{i}}&\mbox{ if $n\geq\hat{n}$}
	\end{cases}
	\end{gathered}
\end{equation}

In this approach, we define a set of scaling factors, $\mathbb{I}(\hat{\lambda}_{i},\hat{n})$, which serve to address both types of non-uniformities. Here, $\hat{\lambda_i}$ represents the irregularity associated with specific RoPE dimensions, while $\hat{n}$ characterizes deviations related to token positions. Concretely, we apply $\mathbb{I}(\hat{\lambda}_{i},\hat{n})$ to adjust the rotational angle assigned to the $i^{th}$ dimension in the RoPE mechanism; the scaling coefficient $\hat{\lambda_i}$ and the position threshold $\hat{n}$ dictate this process. For the first $\hat{n}$-1 tokens, the scaling parameter $\hat{\lambda_{i}}$ is not utilized, thus preserving the default RoPE rotational angle of $\frac{n}{\beta^i}$. However, once token positions reach or exceed $n\ge\hat{n}$, the scaling factor is introduced.

Assuming a desired context window length of $L'$, the task is to determine the set of optimal rescale coefficients ($\mathbb{I}(\hat{\lambda}_{0},\hat{n})$, $\mathbb{I}(\hat{\lambda}_{1},\hat{n})$,...,$\mathbb{I}(\hat{\lambda}_{i},\hat{n})$,...) across all RoPE dimensions from $1$ to $d$. This enables the target $\text{LLM}$, after modifying the $\text{RoPE}$ with these rescale factors, to minimize the next-token prediction loss $\mathcal{L}$ (i.e., the perplexity) for a collection of input sequences $\mathbf{X}$ with length $L'$.

\subsection{Searching the Non-uniform Position Interpolation}

To address the issue outlined in Eq.~\ref{eq:problem}, we propose a straightforward but powerful approach that identifies the optimal RoPE rescale factors, thereby leveraging the diverse non-uniform characteristics present within the multidimensional position embedding.

\textbf{Search space}. We construct an extensive search space to accommodate both forms of non-uniformity. The search space is summarized in Table~\ref{tbl:searchspace}. More precisely, we permit independent rescale factors for each RoPE embedding dimension, enabling tailored adaptation. For simplicity in the design of the search space, the parameters $\lambda_i$ and $\hat{n}$ are directly explored rather than exhaustively searching over $\mathbb{I}(\hat{\lambda}_{i},\hat{n})$, with the understanding that $\hat{\lambda}_{i}=1/\lambda_i$. As indicated in Table~\ref{tbl:searchspace}, the value of $\lambda_i$ is chosen from a range starting at 1.0 (which equates to standard extrapolation) up to $s \times 1.25$ (representing interpolation wider than what is provided by PI), incremented by 0.01. Here, $s$ denotes the extension ratio for the target context window.

The parameter $\hat{n}$ determines how many of the first token positions preserve the original RoPE embedding, foregoing position interpolation. In practice, $\hat{n}$ is chosen from the set \{0, 1, 2, 4, 8, 12, 16, 20, 24, 28, 32, 64, 128, 256\}. If $\hat{n}=0$, every token position relies solely on the optimized rescale factors.

\textbf{Evolution-based search}. The search space outlined in Table~\ref{tbl:searchspace} encompasses a wide array of positional interpolation options, making it challenging to search efficiently. For instance, an extension factor of $s=4\times$ results in $400^{128/2}\times14 = 4\times10^{167}$ possible configurations. As the extension ratio increases, the search space grows at an exponential rate. To tackle this complexity, we employ evolution search~\cite{guo2020single} and present two optimization strategies that substantially improve the efficiency of the search process. The general search methodology is depicted in Algorithm~\ref{alg:search}.

\textit{Optimized initial population generation}. Rather than selecting all $P$ rescale factors in the initial population at random, we ensure that the three RoPE rescale factors associated with PI, NTK, and YaRN are explicitly included as individuals. The remaining $P-3$ individuals are produced by randomly mutating these three rescale factors, each with a mutation probability of $p$.

\textit{Monotonically non-decreasing constraint}. Following the creation of the initial population, we evaluate each candidate by calculating its LLM perplexity. This involves applying the specific RoPE rescale factors to the target LLM, then determining the perplexity for the input $\mathbf{X}$. The top-k candidates, ranked by performance, are selected as parents for the next evolutionary step. Nonetheless, due to the immense search space, simple mutation and crossover can yield suboptimal candidates and trigger redundant perplexity evaluations. This inefficiency is further amplified when $L'$ is large, as each perplexity computation is computationally expensive.

To mitigate this issue, we enforce a monotonic non-decreasing constraint on the sampled RoPE rescaling factors such that $\lambda_i \leq \lambda_{i+1}$ holds. We only consider RoPE configurations adhering to this constraint during their application to LLMs for perplexity assessment, which effectively narrows the search space and lowers the associated computational burden. More concretely, we stipulate that the values of $\lambda_i$ should rise monotonically across the RoPE dimensions (with $i=0,\ldots,63$). The rationale for this monotonicity across dimensions is informed by NTK theory~\cite{jacot2018neural,tancik2020fourier,ntk}, which indicates that lower dimensions, corresponding to higher frequencies, necessitate less interpolation (smaller $\lambda_i$ values), while higher dimensions with lower frequencies are better suited for increased interpolation (larger $\lambda_i$ values).

\begin{algorithm}[t]
	\small
	\caption{The search algorithm}
	\textbf{Input:} target LLM, input samples $\mathbf{X}$,   population size $P$, mutation size $N_1$, crossover size $N_2$, max iterations $\mathcal{T}$, mutate probability $p$\\

	\begin{algorithmic}[1]
		\label{alg:search}
		\STATE Initialize Top-k as the empty set $\phi$;
		\STATE Set $\text{P}_0$ to the result of \textit{Initialize\_population\_with\_optimization} with parameters ($P$, $\mathbf{X}$, $p$);
		\FOR{$i=1$ to $\mathcal{T}$}
		\STATE \textit{Compute\_perplexity} (LLM, $\text{P}_{i-1}$, $\mathbf{X}$);
		\STATE  Update Top-k with \textit{Update\_Topk} (Top-k, $\text{P}_{i-1}$);
		\STATE Form $\text{P}_{mutation}$ by invoking \textit{Mutation\_with\_mono\_constraint} on Top-k using $N_1$ and $p$;
		\STATE Generate $\text{P}_{crossover}$ by applying \textit{Crossover\_with\_mono\_constraint} to Top-k and $N_2$;
		\STATE Set $\text{P}_i$ as the union of $\text{P}_{mutation}$, $\text{P}_{crossover}$, and Top-k;
		\ENDFOR
		\STATE Output the member in Top-k exhibiting the minimal perplexity;
	\end{algorithmic}
\end{algorithm}

\textbf{8$\times$ extension without fine-tuning}. Through evolutionary search, our approach successfully discovers optimal non-uniform RoPE rescale factors, retaining critical positional and dimensional features to reduce information degradation from interpolation. As shown in Fig.\ref{fig:non-ft}, this enables us to expand LLaMA2’s context window from 4k to 32k tokens without any need for fine-tuning. Conversely, alternative strategies like PI as well as non-uniform NTK and YaRN exhibit a sharp rise in perplexity when the context window is increased beyond 2$\times$.

\subsection{Extending LLM Context Window to 2048K}
\label{sec:extendto2048k}

\textbf{Progressive extension to 2048k}. In this section, we present our strategy for expanding the context window of pre-trained LLMs beyond the standard 4k tokens to as much as 2048k. Our approach utilizing non-uniform positional interpolation enables up to an 8$\times$ context window increase with no fine-tuning required. When greater expansion is needed—for example, a 512$\times$ extension—fine-tuning becomes indispensable. One possibility involves identifying appropriate RoPE rescaling coefficients corresponding to the 2048k target length, followed by fine-tuning. Nevertheless, this process is hindered by the immense computational costs involved. Furthermore, our practical attempts reveal that achieving satisfactory fine-tuning outcomes at such large extension ratios is particularly difficult (refer to Appendix).

Fortunately, LongRoPE demonstrates strong performance not only with the original model but also when applied to fine-tuned, longer-context LLMs. As a result, we propose a streamlined, stepwise approach that reaches the desired context length of 2048k in merely 1k fine-tuning iterations, while constraining the training length to within 256k.

$\Diamond$ \textit{LongRoPE search for 256k context extension of pre-trained LLM}. Using LLaMA2 as a case study, we perform a search targeting context window lengths of 128k and 256k. At this level, the model’s context is expanded by factors of 32$\times$ and 64$\times$, respectively.

$\Diamond$ \textit{Fine-tuning to 256k}. Next, the pre-trained LLM is further fine-tuned to accommodate a context window of 256k. The process begins by fine-tuning LLaMA2 for 400 steps with RoPE rescaled factors set to 128k. After completing this phase, the RoPE rescaled factors are updated to 256k on the obtained checkpoint, followed by another 600 fine-tuning steps. This approach is demonstrated to be more effective than immediate fine-tuning to 256k.

$\Diamond$ \textit{Scaling fine-tuned extended LLM to 2048k via LongRoPE search}. In the final step, we apply an additional search phase to the already fine-tuned LLM with a context length of 256k. This process expands the model's context window up to a remarkable 2048k, achieved without requiring additional fine-tuning. Consequently, the overall extension ratio reaches $512\times$.

\textbf{Degradation in performance over shorter context windows}. Expanding the context window length to an extensive 2048k tokens introduces a decline in model effectiveness within the confines of the original context window. This phenomenon, previously identified in literature on positional interpolation~\cite{pi}, occurs because position embeddings for the initial context window are compressed into a significantly smaller interval in the higher-dimensional embedding space. Such compression can harm the language model's ability to process sequences effectively. For example, applying a 512$\times$ extension leads to a substantial congestion of positions within the original 4k token window.

To address this issue, we conduct an additional evolution search on the extended LLM, specifically tuning the RoPE rescale factors for shorter context windows (such as 4k and 8k). The upper limit of the searched $\lambda$ values is lowered, reflecting the diminished need for positional interpolation at these shorter lengths. At inference time, the LLM automatically selects the appropriate RoPE rescale factors for the given context.

 \section{Experiments}

\label{sec:exp}
\subsection{Setup}
\textbf{Evaluation Tasks and models}. LongRoPE is integrated with both LLaMA2-7B and Mistral-7B, and its effectiveness is assessed across three dimensions: (1) the perplexity scores achieved by large language models with extended context lengths when processing lengthy documents; (2) the Passkey retrieval challenge, which evaluates how well the model can identify a specific passkey hidden among a large volume of distracting information; and (3) conventional LLM benchmark tests conducted using a restricted 4096-token context window.

\textbf{Fine-tuning}. For LLaMA2, fine-tuning is performed using a learning rate set to 2e-5 with linear decay and a global batch size of 32. The model is fine-tuned over 400 steps on the Redpajama~\cite{redpajama} dataset, with the data divided into 128k-token chunks framed by BOS and EOS tokens. Subsequently, continuing from the resulting checkpoint, an extra 600 training steps are conducted to expand the context window to 256k tokens. Training on the 128k context window utilizes 8 A100 GPUs within the distributed training infrastructure~\cite{cube}, whereas 16 A100 GPUs are required for the 256k context size. For Mistral, we adopt a fixed learning rate of 1e-6 and a global batch size of 64. Following the YaRN configuration~\cite{yarn}, both 128k and 256k models are each trained for 400 steps using the Together Computer's Long-Data Collections~\cite{mistraldata} at a sequence length of 16k. Training runs on 4 A100 GPUs.

\textbf{Search}. When the target window size does not exceed 256k, we set the parameters as follows: $P$=64, $N_1$=$N_2$=16, $p$=0.3, and $\mathcal{T}$=40, selecting the top-32 candidates for mutation and crossover during every iteration. Perplexity evaluation utilizes 5 randomly selected samples from the PG19 validation dataset, with each sample meeting or exceeding the target context length. For target window sizes above 512k, we reduce the population size, as well as the number of candidates for mutation and crossover, by half. In this case, perplexity is determined on 3 randomly chosen validation samples from Pile-Books3~\cite{pile}.

\textbf{Baselines}. For achieving a 2048k context window, we performed fine-tuning on models using both 128k and 256k sequence lengths. The resulting models are designated as LongRoPE-2048k (ft=128k) and LongRoPE-2048k (ft=256k), corresponding to LLaMA2 and Mistral, respectively. Our evaluation involves comparing these four models against leading context window extension baselines. In particular, we consider open-source LLMs that have been fine-tuned with positional interpolation strategies, including PI, NTK, and YaRN. Among the compared baselines are Together-32k~\cite{together}, Code LLaMA~\cite{codellama}, LongLoRA-full-FT-100k~\cite{longlora}, YaRN-LLaMA, and YaRN-Mistral~\cite{yarn}.

\subsection{Main Results}

\label{sec:mainresult}
\textbf{Long sequence language modeling within 256k}. To start, we benchmark our approach against leading extended LLMs on tasks capped at 256k tokens. To highlight the model's versatility, evaluations are carried out on both the Proof-pile~\cite{pg19} and the PG19~\cite{pile} test sets. Perplexity scores are computed for several context window sizes, using a sliding window method with a stride of 256. For PG19, the analysis encompasses the full test set, which consists of 100 documents. In the case of Proof-pile, our methodology is aligned with YaRN~\cite{yarn}, involving a random selection of 10 passages, each extending to a minimum of 128k tokens.

Table~\ref{tbl:proof-pile} and Table~\ref{tbl:pg19} present a comparison of the perplexity scores for LLaMA2 and Mistral, each extended by various interpolation strategies, on the Proof-pile and PG19 datasets. Two main findings emerge: \textbf{(1)} Our extended models demonstrate a consistent reduction in perplexity as the evaluation sequence length increases from 4k to 256k, indicating enhanced utilization of long-range context. \textbf{(2)} Despite being tested with a context window up to 16$\times$ longer than standard, which often impairs short-context performance, the LongRoPE-2048k variants surpass leading existing models within a 256k context length.

\textbf{Long sequence language modeling beyond 2000k}. To assess performance on exceptionally lengthy texts, we utilize the Books3~\cite{pile} dataset. For practical evaluation, 20 books longer than 2048k tokens are randomly chosen, applying a sliding window approach of 256k.

Table~\ref{tbl:books3} illustrates that LongRoPE is effective in expanding the context window of both LLaMA2-7B and Mistral-7B models to 2048k tokens, while maintaining perplexity levels on par with, or superior to, baseline models at shorter context lengths (8k-128k). Notably, there are substantial performance variations between the 2048k configurations of LLaMA2 and Mistral. Although Mistral surpasses baseline models at shorter context lengths, its perplexity increases to over 7 for contexts longer than 256k. The performance of LLaMA2 is consistent with prior expectations, demonstrating a consistent decrease in perplexity as the context length increases, except for slight increases at 1024k and 2048k. Additionally, for LLaMA2, LongRoPE-2048k yields improved results when fine-tuned with a window of 256k compared to 128k, which can be attributed to the reduced secondary extension ratio (i.e., 8$\times$ rather than 16$\times$). Conversely, Mistral achieves better performance with a fine-tuning window of 128k. This discrepancy is primarily due to the use of a 16k training length in both the 128k and 256k fine-tuning setups for Mistral, as adopted from YaRN, thereby constraining the model’s ability to generalize to longer context windows post fine-tuning.

\textbf{Passkey retrieval}. Next, we investigate how the usable context window influences performance in generative settings. To do so, we adopt a synthetic passkey retrieval task described by~\cite{passkey}, wherein the model must locate a randomly embedded five-digit passkey within an extended document. Details of the prompt structure can be found in the appendix. The experiment involves 10 independent trials, with the passkey randomly situated at uniformly distributed positions throughout the evaluation context length.

Fig.~\ref{fig:passkey} illustrates a comparison of retrieval accuracy across various baseline models. The performance of previous models declines sharply, with accuracy falling to zero past 128k tokens. In stark contrast, our LongRoPE-LLaMA2-2048k (ft=256k) demonstrates robust retrieval capabilities even in the difficult setting of passkey identification among millions of tokens, maintaining a high accuracy rate ($\ge$90\%) consistently from 4k up to 2048k tokens. Similarly, LongRoPE-Mistral-2048k (ft=128k) sustains perfect accuracy (100\%) until approximately 1800k tokens, after which accuracy decreases to 60\% at 2048k. This result is consistent with the patterns observed in Table~\ref{tbl:books3}, where there is a modest increase in perplexity at the 2048k token length.

\textbf{Evaluation on standard benchmarks within the initial context window}. We assess the performance of LongRoPE-2048k models on their original context window via the Hugging Face Open LLM Leaderboard~\cite{huggingfaceboard} under both zero-shot and few-shot paradigms. Specifically, the evaluation leverages 25-shot ARC-Challenge~\cite{allenai:arc}, 10-shot HellaSwag~\cite{zellers2019hellaswag}, 5-shot MMLU~\cite{mmlu}, and 0-shot TruthfulQA~\cite{lin2021truthfulqa}.

As indicated in Table~\ref{tbl:commonsense}, our models perform similarly to previous systems on the standard benchmark tailored for limited context, and surpass the original Mistral by 0.5\% on TruthfulQA. While LongRoPE-LLaMA2-2048k, after being fine-tuned with a 256k context, exhibits a marginally greater decline in performance, it generally stays within acceptable bounds across the majority of tasks.

\subsection{Ablation Results}

\textbf{Effectiveness of the second positional interpolation}. Within our progressive extension strategy, we apply a second round of non-uniform positional interpolation, utilizing our search algorithm, to fine-tuned extended LLMs. To assess the efficacy of this approach, we conduct experiments using our fine-tuned LLaMA2-256k model. We scale the model to 512k, 1024k, and 2048k contexts employing PI and YaRN for comparison. The results in Table~\ref{tbl:secondsearch} indicate that our non-uniform positional interpolation maintains perplexity at a steady level. In comparison, perplexity with PI and YaRN escalates rapidly as the extension ratio increases.

\textbf{Recovery performance at reduced context windows.} To address the degradation in model performance at shorter context ranges, we optimize the RoPE parameters of LongRoPE-2048k using our search procedure. In particular, we limit the maximum scaling factors during the search phase, thereby discouraging excessive interpolation for short contexts of 4k and 8k tokens. Table~\ref{tbl:shortperformance} provides perplexity results for LongRoPE-LLaMA2-2048k evaluated on Proof-pile at these shorter context lengths, as well as corresponding mean scores on standard LLM benchmarks. The data indicate marked improvements in performance at reduced context lengths.

\textbf{Analysis on the two forms of non-uniformities}. To further understand the impact of the two sources of non-uniformity, we conduct ablation studies to isolate their individual contributions to overall performance. Two sets of experiments are designed: (i) for LLaMA2-7B, we extend the context window to 16k and 32k tokens employing different approaches—Position Interpolation (PI), searching solely for the optimal RoPE dimension, and jointly searching for both non-uniformities; (ii) for our 256k-length fine-tuned LLaMA2, we scale up to a 2048k context window using a similar process. Perplexity measurements are obtained without additional fine-tuning. Table~\ref{tbl:searchdimension} indicates that introducing non-uniformity in the RoPE dimension yields considerable perplexity reduction compared to standard linear interpolation from PI. Non-uniform treatment of token positions delivers noticeable improvements at sequence lengths of 16k and 32k, but its effect is diminished at 2048k, perhaps owing to the challenges at such long sequence lengths. Notably, simply retaining the original tokens without any form of interpolation ceases to offer benefits; investigating this phenomenon in greater depth is left for future studies.

\section{Related Works}

Beyond position interpolation techniques, this section reviews alternative strategies found in related literature.

\textbf{Retrieval-based} approaches leverage external memory components to store distant contextual information and employ retrieval mechanisms to access pertinent documents during inference~\cite{longllama,wang2023augmenting,borgeaud2022improving}. Implementing these methods generally requires substantial architectural changes to the underlying LLMs. By contrast, our method remains comparatively lightweight, necessitating only modest adjustments to positional embeddings. Furthermore, it accommodates a broader array of long context applications aside from retrieval, including tasks such as extensive document summarization and few-shot learning.

\textbf{Attention-based context window extensions}. In addition to positional embedding interpolation techniques, certain studies have explored extending the input context by adjusting the attention mechanism within the standard LLM context window length~\cite{han2023lminfinite, streamingllm,parallelwindow}. The principal concept involves addressing the problem of attention explosion at newly introduced positions through the use of innovative attention masks. Such methods serve as a complement to approaches based on positional interpolation.

\textbf{Fine-tuning based approaches} are concerned with optimizing the adaptation of pre-trained LLMs by altering position embeddings to accommodate extended contexts. For instance, methods such as Code LLaMA~\cite{codellama}, LLaMA2 Long~\cite{xiong2023effective}, and ScaledRoPE~\cite{liu2023scaling} adopt an increased base value for RoPE and then fine-tune models on the desired sequence lengths. Our approach, in contrast, enables adaptation across a range of target lengths and supports scaling to beyond 2M tokens. As extending context windows beyond 128k tokens requires significant computational resources, recent techniques such as LongLoRA~\cite{longlora} and PoSE~\cite{zhu2023pose} aim to reduce GPU memory and compute demands associated with long-context fine-tuning. Notably, our method complements these efficiency-oriented approaches.

\section{Conclusion}

In this study, we introduce LongRoPE, a technique that significantly increases the context window of LLMs to an extraordinary 2048k tokens, all while preserving their performance on shorter, original context lengths. Our approach leverages two distinct types of non-uniformities within RoPE positional embeddings, utilizing an efficient evolutionary search framework. This approach yields two major advantages: it offers effective initialization for subsequent fine-tuning, and facilitates an 8$\times$ increase in context window size even in the absence of fine-tuning. Expanding upon this, we introduce a stepwise extension methodology, where LLMs fine-tuned on a 256k context are further extended to a 2048k context window without the need for additional fine-tuning. Comprehensive experimental results demonstrate the efficacy of LongRoPE. We anticipate that our LongRoPE-2048k models will open up a range of novel applications requiring long context and encourage further investigations in this field.

\section*{Broader Impacts} The research described in this paper is aimed at contributing to the progress of Machine Learning as a discipline. While our findings may have various societal implications, we do not identify any particular issues that require emphasis in this context.

	\balance

\end{document}